\subchapter{Exploring the secure world}{Objective: 
	Observe Exception Levels and the secure world, build secure world software.}

\section{Observing Exception Levels}

In this section, we are going to patch secure world software to display some information.
The software we are going to patch is:
\begin{itemize}
  \item ARM Trusted Firmware (TF-A)
  \item the Trusted Execution Environment: OP-TEE
\end{itemize}

\subsection{Preparing to patch}

\begin{bashinput}
$ git clone https://review.trustedfirmware.org/TF-A/trusted-firmware-a.git --revision 569e16caad976a0684147da1ecc6333fd9b7f813
$ git clone https://github.com/OP-TEE/optee_os.git --revision 18b424c23aa5a798dfe2e4d20b4bde3919dc4e99 
\end{bashinput}

Once you are done modifying the code in the next sections, do

\begin{bashinput}
$ git commit -s -m 'Bootlin Security Training'
$ git format-patch HEAD~1
\end{bashinput}

This should generate \code{0001-Bootlin-Security-Training.patch}.
Make sure the filename is correct!

Once you are happy with your patch, you can copy it to the appropriate destination to be applied.

For TF-A:
\begin{bashinput}
$ cp 0001-Bootlin-Security-Training.patch $HOME/security-imx93-frdm-labs/board_setup/meta-security-training/recipes-bsp/trusted-firmware-a/trusted-firmware-a/
\end{bashinput}

For OP-TEE:
\begin{bashinput}
$ cp 0001-Bootlin-Security-Training.patch $HOME/security-imx93-frdm-labs/board_setup/meta-security-training/recipes-security/optee/optee-os/
\end{bashinput}

And re-build the image:
\begin{bashinput}
$ cd board_setup
$ . openembedded-core/oe-init-build-env
$ bitbake security-training-image
\end{bashinput}

\subsection{Adding logs to TF-A}

We are interested in observing two characteristics of the system:
\begin{itemize}
  \item The current Exception Level
  \item The Security state
\end{itemize}
Once you've determined where that information is stored (e.g. by looking at
\href{}{the ARM documentation}), you can either implement a routine for looking it up
yourself, or by searching for relevant functionality in the TF-A codebase.

\subsection{Adding logs to OP-TEE}

The \href{https://support.arm.com/documentation/ddi0595/2021-06/AArch64-Registers/SCR-EL3--Secure-Configuration-Register}{ARM specification}
states that accessing the Security State from OP-TEE (in fact any
exception level apart from EL3) is "UNDEFINED". So here, let's add some logs to
print the current Exception Level only.

OP-TEE does some serial initialization in
\code{core/arch/arm/kernel/entry_a64.S}, so add a branch instruction there to
our own custom C function right after the console is initialized.

This function should be created next to where \code{console_init} is
defined. Don't forget to also declare this function in the relevant header!

Since OP-TEE does not have a helper for getting the current Exception Level, we
need to replicate what TF-A is doing. Back in TF-A, have a look at the
definition of \code{get_current_el()} and the following macro statement in
\code{include/arch/aarch64/arch_helpers.h}:

\begin{verbatim}
DEFINE_SYSREG_READ_FUNC(CurrentEl)
\end{verbatim}

What is this doing? Replicate this in OP-TEE. The
\href{https://gcc.gnu.org/onlinedocs/gcc/Extended-Asm.html}{Extended ASM GNU
documentation} may be helpful here.

Using the \code{IMG()} macro, show the Exception Level value.

\section{Writing a small Trusted Application (TA)}

In order to understand how Trusted Applications work, we are going to write and deploy one.
More specifically, we are going to write a 
\href{https://optee.readthedocs.io/en/latest/architecture/trusted_applications.html\#user-mode-trusted-applications}{User Mode TA}
which will be loaded from the \href{https://optee.readthedocs.io/en/latest/architecture/trusted_applications.html\#ree-filesystem-ta}{REE filesystem}.
The goal of this Trusted Application will be to try to show the current
Exception Level, just like we've done from OP-TEE.

The requirements for Trusted Applications are documented in the
\href{https://optee.readthedocs.io/en/latest/building/trusted_applications.html}{build} section of the OP-TEE documentation.

We are going to use OP-TEE's TA-devkit to build the TA, so we need to meet those requirements.
In the labs data, under \code{board_setup/meta-security-training/recipes-security/test_ta/files},
the basic infrastructure for building the TA is provided:

\begin{itemize}
  \item The \code{ta/} directory contains the code of the TA.
  \item The \code{host/} directory contains an app to run on the target to use
    the TA.
\end{itemize}

REE FS (Usermode) TAs, contrary to pseudo TAs, are loaded into OP-TEE when the REE attempts to call them.
Therefore we need a client on the REE side to call our TA.
It will need to invoke the command implemented in the previous step.

Go over each \code{TODO} comment in the different files and fill in the missing
pieces.

The host client code is already implemented in
\code{board_setup/meta-security-training/recipes-security/test_ta/files/host},
have a look at the code.

Once you are done, you can run
\begin{bashinput}
$ cd board_setup
$ . openembedded-core/oe-init-build-env
$ bitbake test-ta
\end{bashinput}

The TA will be deployed in \code{board_setup/build/tmp-glibc/deploy/images/freiheit93/test-ta}.

\section{Deploying the TA}

OP-TEE's REE FS Trusted Applications should be installed in \code{/lib/optee_armtz}.
Each one is stored as \code{<UUID>.ta}, we'll deploy ours in the same way:

\begin{bashinput}
$ cp <UUID>.ta /lib/optee_armtz/
\end{bashinput}

Confirm that \code{tee-supplicant} is running, as that's the daemon that will load the TA:
\begin{bashinput}
$ ps | grep tee
\end{bashinput}

You can now trigger your TA:
\begin{bashinput}
$ ./test_ta_host
\end{bashinput}

You should get an exception from OP-TEE, can you guess why?
